% generated by scripts/statutory-qa/make_benchmark_example.py; do not edit
\small
\paragraph{Public input.} Question M027-Q01, paper 4, October 2024 sitting; reference date 2024-10; enrichment not applied.
\begin{quote}\small
\textit{Background.} Kanthi Peiris is a Notary practicing in Colombo. Shantha Ratnayake is her brother-in-law and is the owner of an apartment in Negambo. Shantha wishes to lease out his apartment for 3 years at a monthly rental of Rs.120,000/= to one Ananda Gamage who is willing to take on lease in accordance with the said terms. However, Shantha who lives in Negambo is unable to come to Colombo to sign the Agreement.

\textit{Question.} How can this Agreement be signed by both parties and what steps should Kanthi follow with regard to the registration of this Agreement?
\end{quote}
\paragraph{Proposed gold (private view).} Authority requirement \textsc{statute only}; legal hop count 4; gold Acts 1-1907, 23-1927, 7-1840; status \textsc{proposed gold}, lawyer validation \textsc{pending}.
\begin{enumerate}\setlength{\itemsep}{0pt}
\item Whether a three-year lease of an apartment must be executed notarially, and in whose presence the lessor and lessee must sign.
\item How the agreement can be validly signed when the lessor cannot attend before the Colombo notary - execution before more than one notary at different times and places, or execution by a duly authorised attorney under a registered power of attorney.
\item What registration steps the notary must take for a lease of land situated in a district other than the one in which she resides and practises, and within what time.
\end{enumerate}
\begin{tabular}{@{}p{0.30\linewidth}p{0.66\linewidth}@{}}
\toprule Indispensable provision & Excerpt checked verbatim against the corpus \\ \midrule
Prevention of Frauds Ordinance No. 7 of 1840, s.\,2(1)(a) & shall be in force or avail in law unless - (a) the relevant deed or instrument shall be in writing, signed by every executant or by any person duly authorised by such executant and the witnesses in the presence of a lic… \\
Notaries Ordinance No. 1 of 1907, s.\,31 rule (12) & (12) He shall not authenticate or attest any deed or instrument unless the person executing the same and the witnesses shall have signed the same in his presence and in the presence of one another, and unless he shall h… \\
Notaries Ordinance No. 1 of 1907, s.\,31 rule (22) & (22) He shall not authenticate or attest any deed or instrument in any area other than that in which he is authorized to practice, nor in any language other than that in which he is authorized to practice nor authentica… \\
Notaries Ordinance No. 1 of 1907, s.\,31 rule (25) & (25) Where any deed or instrument is executed or acknowledged before more than one notary- (a) the notary who first attests such deed or instrument shall comply with all the requirements of rule (20), and every other no… \\
Notaries Ordinance No. 1 of 1907, s.\,32(2)(i) & (2) In the case of any deed or instrument which is to be executed by two or more parties, both or all of whom, as the case may be, do not sign the deed or instrument at the same time and place- (i) the deed or instrumen… \\
Notaries Ordinance No. 1 of 1907, s.\,31 rule (30A) & (30A) It shall be the duty of every notary to submit for registration to the Registrar, every deed or instrument attested by him before the expiry of thirty days from the date of attestation thereof: Provided that, wher… \\
Registration of Documents Ordinance No. 23 of 1927, s.\,8(b) & (b) if executed or made after the commencement of this Ordinance, all instruments, including wills, decrees and orders of any court or authority, and awards, which purport or operate to create, confer, declare, limit, a… \\
Registration of Documents Ordinance No. 23 of 1927, s.\,14(1) & (1) Every instrument presented for registration shall be registered in the book allotted to the division in which the land affected by the instrument is situated \\
Notaries Ordinance No. 1 of 1907, s.\,2 & Every appointment to the office of notary shall be by warrant granted by the Minister, and shall specify the area within which, and the language or languages in which, the person appointed is authorized to practice. \\
\bottomrule\end{tabular}

\noindent Supporting (not scored): 1-1907/s37, 23-1927/s26, 23-1927/s7, 4-1902/s2, 4-1902/s3, 4-1902/s3A, 4-1902/s3B.

\noindent Reasoning edges (source $\to$ target):

\noindent\begin{tabular}{@{}lp{0.74\linewidth}@{}}
\textsc{jointly required} & 7-1840/s2 $\to$ 1-1907/s31; 23-1927/s8 $\to$ 23-1927/s14 \\
\textsc{qualifies} & 1-1907/s31 $\to$ 1-1907/s32; 23-1927/s8 $\to$ 23-1927/s7 \\
\textsc{procedurally requires} & 7-1840/s2 $\to$ 1-1907/s31; 1-1907/s31 $\to$ 23-1927/s14 \\
\textsc{cross references} & 1-1907/s31 $\to$ 4-1902/s2; 4-1902/s3A $\to$ 7-1840/s2 \\
\textsc{defines} & 1-1907/s2 $\to$ 1-1907/s31 \\
\end{tabular}
\paragraph{Draft answer (agent-written, unvalidated).}
\textit{Issue.} How a three-year lease of a Negombo apartment at Rs.120,000 a month can be signed by a lessor who cannot come to Colombo and by a lessee who can, and what a Colombo notary must then do to register the instrument.

\textit{Rule.} A lease of immovable property for a period exceeding one month is of no force in law unless it is in writing, signed by every executant, or by a person duly authorised by the executant, and by the witnesses in the presence of a licenced notary present at the same time and in the presence of one another, and attested by that notary (Prevention of Frauds Ordinance s.2(1)(a)). The notary may not attest unless the executant and the witnesses signed in her presence and in the presence of one another (Notaries Ordinance s.31 rule (12)), and she may not attest in any area other than the area in which she is authorised to practise (s.31 rule (22)). The Ordinance nevertheless contemplates a deed executed or acknowledged before more than one notary and divides the attestation duties between the first attesting notary and every other attesting notary (s.31 rule (25)), and s.32(2) provides for parties who do not sign at the same time and place, deeming the instrument executed for the purposes of rules (6), (7), (23) and (25) when it is first signed. Alternatively the lessor may authorise an attorney: a power of attorney executed in Sri Lanka must be executed before two witnesses and attested by a notary (Powers of Attorney Ordinance s.2), must meet the additional requirements of s.3A where the transaction falls within s.2 of the Prevention of Frauds Ordinance, must be submitted for registration with the Registrar-General within one month (s.3(1)(b)), and the attesting notary must check the Registrar-General's volumes and folios and state in her attestation that it has not been revoked (s.3B), keep a true copy with her protocol and forward a copy with the duplicate (Notaries Ordinance s.31 rule (30)), and take the attorney's affidavit that the power is genuine, in force and the grantor alive. As to registration, an instrument creating an interest in land is deemed to affect land and so is registrable (Registration of Documents Ordinance s.8(b)), an unregistered instrument is void against a subsequent duly registered instrument taken for value (s.7(1)), registration must be effected in the book allotted to the division in which the land is situated (s.14(1)), and the notary acting for a party may present it (s.26(1)(d)). The notary must submit the deed for registration within thirty days of attestation, or sixty days where it is to be registered outside the jurisdiction in which she practises (Notaries Ordinance s.31 rule (30A)), send the duplicate and Form F1 list to the Registrar of Lands of her own district by the fifteenth of the following month (rule (26)), send a certified copy to the Registrar of Lands of the district where the land lies (rule (29)), and use due diligence in registering where she has instructions and funds (s.37).

\textit{Application.} The lease is for three years, well outside the one-month exclusion, so notarial execution is compulsory and Shantha cannot simply sign at home and send the paper to Colombo. Kanthi practises in Colombo and the apartment and Shantha are in Negombo, so rule (22) prevents her attesting Shantha's signature outside the area she is authorised for. Two lawful routes remain. First, execution before more than one notary under rule (25): Ananda Gamage signs the agreement before Kanthi in Colombo with two witnesses, and Shantha signs the same instrument before a notary authorised to practise where he is, each notary attesting the signatures taken before her or him, each numbering the instrument in his or her own protocol register, and s.32(2) accommodating the fact that the parties sign at different times and places. Second, Shantha grants a power of attorney, attested by a notary before two witnesses and complying with s.3A because the lease falls within s.2 of the Prevention of Frauds Ordinance, registers it with the Registrar-General within one month, and his attorney then signs before Kanthi in Colombo along with Ananda; Kanthi must satisfy herself from the Registrar-General's volumes and folios that the power stands unrevoked and say so in her attestation, take the attorney's affidavit, and keep and forward copies of the power. Either way she must have an assurance that the stamp duty is provided before attesting. On registration, the lease creates an interest in land and is registrable; it must go to the land registry for the division in which the Negombo apartment is situated, not Colombo. Because that registration is outside the jurisdiction in which Kanthi practises, she has sixty days from attestation to submit it, she must transmit the duplicate and Form F1 list to the Registrar of Lands of her own district by the fifteenth of the following month, and she must send a certified copy to the Registrar of Lands of the district in which the apartment lies. If the apartment forms part of a registered Condominium Property, she should also check whether the interest has to be registered under the Registration of Title Act rather than in the ordinary land register. The scenario does not state Shantha's judicial zone, whether the apartment is registered condominium property, or whether the parties have instructed and funded Kanthi to register, so those points are stated conditionally.

\textit{Conclusion.} The agreement can be signed either before two notaries at different times and places, Ananda before Kanthi in Colombo and Shantha before a notary where he is, or by an attorney holding a duly attested and registered power of attorney signing for Shantha before Kanthi. Kanthi must then have the duty stamped, submit the instrument for registration in the land registry division covering the Negombo apartment within sixty days of attestation, transmit the duplicate with the Form F1 list to the Registrar of Lands of her own district by the fifteenth of the following month, and send a certified copy to the Registrar of Lands of the district where the apartment is situated.

\noindent Claim-level provenance, first two of 16:
\begin{itemize}\setlength{\itemsep}{0pt}
\item [rule] A lease of immovable property for a period exceeding one month is of no force in law unless it is in writing, signed by every executant or a person duly authorised by the executant and by the witnesses in the presence of a licenced notary, and attested by that notary. \hfill $\leftarrow$ 7-1840/s2, 1-1907/s31, 7-1840/s2
\item [rule] A notary may not attest a deed unless the executant and the witnesses signed it in her presence and in the presence of one another. \hfill $\leftarrow$ 1-1907/s31
\end{itemize}
